\documentclass[10pt]{article}
\usepackage[T1]{fontenc}
\usepackage[utf8]{inputenc}
\usepackage[french]{babel}
\usepackage{amsmath,amssymb,mathrsfs,mathpazo}
\usepackage{geometry}
\geometry{a4paper,margin=2cm}
\usepackage{graphicx,array,tabularx,tabularray,colortbl,multirow,multicol,booktabs,enumitem,pifont,ulem,fancybox,diagbox,cancel}
\usepackage{pstricks,pstricks-add,pst-plot,pst-tree,pst-node,pst-text,pst-func,pst-math,pst-bspline,pst-3dplot,pst-eucl,pst-solides3d}
\let\pscalculate\undefined
\usepackage{pst-bezier}
\usepackage{tikz,pgfplots,tkz-tab}
\usetikzlibrary{arrows,arrows.meta,patterns,calc,positioning,trees,decorations.markings,shapes.geometric,shapes.misc,shapes.arrows,intersections,decorations.pathmorphing,decorations.pathreplacing,backgrounds,shadows,babel}
\pgfplotsset{compat=1.18}
\usepgfplotslibrary{fillbetween}
\usepackage[np]{numprint}
\usepackage{listings,makecell,pas-tableur,twoopt}
\let\barre\undefined
\usepackage[pstricks]{tablvar}
\definecolor{darkgreen}{RGB}{0,100,0}
\providecommand{\red}{\color{red}}
\providecommand{\blue}{\color{blue}}
\providecommand{\green}{\color{green}}
\providecommand{\black}{\color{black}}
\providecommand{\white}{\color{white}}
\providecommand{\gray}{\color{gray}}
\providecommand{\cyan}{\color{cyan}}
\providecommand{\magenta}{\color{magenta}}
\providecommand{\yellow}{\color{yellow}}
\providecommand{\violet}{\color{violet}}
\providecommand{\orange}{\color{orange}}
\providecommand{\euro}{\texteuro}
\providecommand{\EUR}{\texteuro}
\providecommand{\ieme}{\textsuperscript{e}}
\providecommand{\degres}{\ensuremath{^\circ}}
\providecommand{\degree}{\ensuremath{^\circ}}
\providecommand{\vv}[1]{\overrightarrow{#1}}
\providecommand{\Coord}[1]{#1}
\providecommand{\decosix}{}
\providecommand{\starredbullet}{\ensuremath{\bullet}}
\pagestyle{empty}
\setlength{\parindent}{0pt}
\usepackage{adjustbox}\DeclareUnicodeCharacter{2236}{:}\DeclareUnicodeCharacter{2190}{\ensuremath{\leftarrow}}
\providecommand{\R}{\mathbb{R}}
\providecommand{\N}{\mathbb{N}}
\providecommand{\Z}{\mathbb{Z}}
\providecommand{\Q}{\mathbb{Q}}
\providecommand{\D}{\mathbb{D}}
\providecommand{\C}{\mathbb{C}}
\providecommand{\e}{{\mathrm{e}}}
\providecommand{\vect}[1]{\overrightarrow{#1}}
\providecommand{\vectt}[1]{\overrightarrow{\mathrm{#1}}}
\providecommand{\vvt}[1]{\overrightarrow{\mathrm{#1}}}
\providecommand{\np}[2][]{{#2~#1}}
\providecommand{\vv}[1]{{\overrightarrow{#1}}}
\renewcommand{\d}{{\mathrm{d}}}
\providecommand{\degres}{{^\circ}}
\providecommand{\degre}{{^\circ}}
\providecommand{\blue}{\color{blue}}
\providecommand{\red}{\color{red}}
\providecommand{\enskip}{\quad}
\providecommand{\llbracket}{\lbrack\!\lbrack}
\providecommand{\rrbracket}{\rbrack\!\rbrack}
\providecommand{\euro}{{\text{€}}}
\providecommand{\no}{{\text{n°}}}
\providecommand{\interval}[2]{{[#1;#2]}}
\providecommand{\intervFO}[2]{{[#1;#2[}}
\providecommand{\dotfill}{\ldots}

\providecommand{\ds}{\displaystyle}
\providecommand{\textding}[1]{#1}
\providecommand{\Oij}{\(\mathrm{O};\vec{\imath},\vec{\jmath}\)}
\providecommand{\Oijk}{\(\mathrm{O};\vec{\imath},\vec{\jmath},\vec{k}\)}
\providecommand{\Ouv}{\(\mathrm{O};\vec{u},\vec{v}\)}
\providecommand{\intervFF}[2]{\ensuremath{[#1;#2]}}
\renewcommand{\euro}{\ifmmode\text{\texteuro}\else\texteuro\fi}
\DeclareUnicodeCharacter{00B0}{\ifmmode^{\circ}\else\textdegree\fi}
\DeclareUnicodeCharacter{20AC}{\ifmmode\text{\texteuro}\else\texteuro\fi}
% Énoncés : Éducation nationale. Transcription et corrigés : APMEP, Denis Vergès et contributeurs. Découpage et classement : CoopMaths, Sébastien Lozano, Éric Elter et contributeurs. Adaptation de la présentation : Maths Explorer.
% CC BY-SA : https://coopmaths.fr/www/about/
% https://coopmaths.fr/alea/static/eam/2026/tex/eam_specialite_2026_04_polynesie_2_cor.tex
\begin{document}
\noindent\textit{APMEP / CoopMaths. CC BY-SA. Présentation adaptée par Maths Explorer.}\par\medskip
\emph{Pour $t \in [0~;~13]$, $f(t) = \left(t^2 + t + 1\right)\e^{-t}$ où $t$ est le nombre
d'heures écoulées depuis l'injection et $f(t)$ la concentration du médicament dans le sang,
en g/L.}

\begin{enumerate}
	\item $f(0) = \left(0^2 + 0 + 1\right)\e^{0} = 1 \times 1 = 1$.

	Au moment de l'injection, la concentration du médicament dans le sang est de
	1~g/L.

	\item
	\begin{enumerate}
		\item $f$ est un produit $uv$ avec $u(t) = t^2 + t + 1$, $u'(t) = 2t + 1$ et
		$v(t) = \e^{-t}$, $v'(t) = -\e^{-t}$. Pour tout $t$ de $[0 ; 13]$ :
		\[f'(t) = (2t + 1)\e^{-t} + \left(t^2 + t + 1\right)\times\left(-\e^{-t}\right)
		= \e^{-t}\left[(2t + 1) - \left(t^2 + t + 1\right)\right]
		= \e^{-t}\left(-t^2 + t\right)\]
		donc $f'(t) = t(-t + 1)\e^{-t}$.

		\item Le facteur $t$ s'annule en 0 ;\\

		 $-t + 1 \geqslant 0 \iff t \leqslant 1$.

		\begin{center}
			\begin{tikzpicture}
				\tkzTabInit[lgt=3.4,espcl=2.6]
				{$t$/1,$t$/1,$-t+1$/1,$t(-t+1)$/1}{$0$,$1$,$13$}
				\tkzTabLine{z,+,,+,}
				\tkzTabLine{,+,z,-,}
				\tkzTabLine{z,+,z,-,}
			\end{tikzpicture}
		\end{center}

		\item Pour tout réel $t$, $\e^{-t} > 0$ : le signe de $f'(t)$ est celui de
		$t(-t + 1)$.

		\begin{center}
			\begin{tikzpicture}
				\tkzTabInit[lgt=3.4,espcl=2.8]
				{$t$/1,$f'(t)$/1,$f$/2}{$0$,$1$,$13$}
				\tkzTabLine{z,+,z,-,}
				\tkzTabVar{-/,+/,-/}
			\end{tikzpicture}
		\end{center}

Il n'est pas demandé de calculer les images.
	\end{enumerate}

	\item Le médicament reste actif tant que $f(t) \geqslant 0,1$. Comme $f$ est décroissante sur $[1~;~13]$ :

\begin{center}
\begin{tblr}{hlines, vlines, colspec={c*{11}{X[c]}}, cells={font=\small}}
$t$ & 3 & 4 & 5 & 6 & 7 & 8 & 9 & 10 & 11 & 12 & 13 \\
$f(t)$ & 0,647 & 0,385 & 0,209 & \textbf{0,107} & 0,052 & 0,025 & 0,011 & 0,005 & 0,002 & 0,001 & 0 \\
\end{tblr}
\end{center}

	$f(6) \approx 0,107 \geqslant 0,1$ et $f(7) \approx 0,052 < 0,1$.

	Le nombre maximal d'heures pendant lequel le médicament reste actif est
	$t_0 = 6$. La nouvelle injection doit être pratiquée au bout de
	6 heures.
\end{enumerate}

\end{document}
